% !TEX root = main.tex
\section{Kết luận}\label{sec:conclusion}

\subsection{Tóm tắt kết quả chính}\label{sec:conclusion:summary}

Đề tài đã xây dựng một khung thực nghiệm hoàn chỉnh cho bài toán tìm đường trong mạng giao thông động trên mạng đường
OpenStreetMap của quận Hà Đông (\NNodes{} nút, \NEdges{} cạnh, \TotalKm{} km), với vận tốc 24 khung giờ được mô phỏng và môi
trường thực có nhiễu tương quan cùng sự cố. Trên \NTest{} truy vấn kiểm tra, các kết quả chính là:
\begin{enumerate}[label=(\arabic*)]
    \item Tính đến sự phụ thuộc thời gian bằng TD-A* giảm regret từ \DijRegret\% (tuyến ngắn nhất) xuống \PoRegret\% (RQ1).
    \item Tái lập kế hoạch đưa regret xuống \RpRegret\%, và replan kết hợp dự báo GCN đạt \RpGnnRegret\% --- tốt nhất trong mọi
    phương pháp khả thi. Replan lười giữ nguyên chất lượng với tốc độ nhanh gấp \LazySpeedup{} lần (RQ2).
    \item Heuristic cận dưới chính xác giảm \LbReduction\% số nút mở rộng so với TD-Dijkstra (heuristic chim bay chỉ giảm
    \TdAstarReduction\%); TD-A* biên dịch Numba nhanh gấp \NumbaSpeedup{} lần bản Python.
    \item RL chọn từng cạnh với quan sát cục bộ chỉ tới đích \RlSuccess\%; quan sát cây kết hợp GNN và học bắt chước nâng lên
    \RlTreeBcSuccess\% nhưng vẫn kém plan-once rõ rệt. RL lai luôn tới đích với regret \RlHybRegret\%, ngang nhưng không vượt
    replan (RQ3).
    \item Định tuyến P90 và Q-learning chọn giờ xuất phát không mang lại lợi ích trong mô phỏng; khi \MaVeh{} xe cùng làm theo
    một chỉ dẫn, thời gian đi tăng \MaNaiveInc\% và phối hợp tuần tự giảm được \MaCoordGain\% (RQ4).
\end{enumerate}

\subsection{Đóng góp của đề tài}\label{sec:conclusion:contrib}

\begin{itemize}
    \item Một khung thực nghiệm mã nguồn mở, tái lập được, chạy trên CPU, từ dữ liệu OpenStreetMap tới báo cáo, trong đó mọi
    con số được sinh tự động từ kết quả chạy.
    \item Một mô hình vận tốc 24 khung giờ có cấu trúc cho Hà Đông, bảo đảm tính FIFO có chứng minh (Bổ đề~\ref{bd:fifofix},
    Mệnh đề~\ref{md:fifoize}).
    \item Một bộ thuật toán tìm kiếm, tái lập kế hoạch, dự báo và định tuyến theo rủi ro được cài đặt hiệu quả và kiểm chứng
    bằng 55 kiểm thử đơn vị.
    \item Một so sánh công bằng giữa tìm kiếm có mô hình và sáu hướng cải tiến học tăng cường trên cùng môi trường, cùng tập
    kiểm tra, cùng chỉ số, kèm phân tích trung thực về khi nào mỗi hướng có và không có lợi.
    \item Ứng dụng minh hoạ trực quan và gói Docker để chạy lại trên máy khác.
\end{itemize}

\subsection{Hạn chế và hướng phát triển}\label{sec:conclusion:future}

\begin{itemize}
    \item \textbf{Dữ liệu tự mô phỏng.} Đây là hạn chế lớn nhất: kết luận định lượng chỉ đúng trong mô phỏng. Bước tiếp theo
    quan trọng nhất là hiệu chỉnh tệp điểm nóng, profile theo giờ, độ nhạy loại đường và mô hình nhiễu bằng dữ liệu GPS, dữ
    liệu xe công nghệ, camera hoặc dịch vụ giao thông thời gian thực.
    \item \textbf{Mạng đường.} Chưa mô hình hoá đèn tín hiệu, hạn chế rẽ trái/quay đầu và xe máy; phạm vi chỉ là quận cũ cộng
    500~m. Khi mở rộng ra toàn thành phố, cần các kỹ thuật tiền xử lý như Contraction Hierarchies phiên bản phụ thuộc thời gian.
    \item \textbf{Dự báo.} GCN hiện chỉ dùng quan sát tức thời; các mô hình không gian--thời gian như DCRNN~\cite{li2018dcrnn}
    hay STGCN~\cite{yu2018stgcn} với lịch sử nhiều bước có thể dự báo tốt hơn.
    \item \textbf{Học tăng cường.} Tinh chỉnh sau học bắt chước nên dùng DAgger đầy đủ hoặc ràng buộc khoảng cách KL với chính
    sách bắt chước để tránh suy giảm; GNN trên đồ thị lân cận thực (thay cho cây) và ngân sách huấn luyện dài hơn; ứng dụng RL
    ở cấp quyết định cao hơn (chọn hành lang, chọn giờ với đặc trưng phong phú hơn).
    \item \textbf{Nhiều xe.} Hàm BPR là mô hình đơn giản; cần mô phỏng vi mô hoặc gán luồng động và tính cân bằng người dùng.
    \item \textbf{Rủi ro.} Dùng ngân hàng kịch bản lớn hơn, mô hình sự cố có xác suất theo từng đoạn đường, hoặc mục tiêu
    trung bình--phương sai thay cho phân vị.
\end{itemize}
